\subsubsection{Lecteurs enrichis et mémoire de fibre}
\label{act21:sec:ch-memoria-fibra}

La différence entre histoire et valeur admet un critère opérationnel. À un
niveau fini \(H_k\), soient \(q_k:H_k\to Z_k\) un lecteur et
\(M_k:H_k\to\mathcal M_k\) le registre conservé. Le lecteur enrichi
est \(\widehat q_k=(q_k,M_k)\). La mémoire distingue des états d'une même
fibre du lecteur ; sa suffisance se vérifie dans le domaine considéré.
Cette formulation intègre à l'article lui-même le développement complet de
l'information de fibre utilisée ci-après.

\begin{proposition}[Récupérabilité à un niveau fini]
\label{act21:prop:ch-recuperabilidad-fibra}
Il existe une inverse de \(\widehat q_k\) sur son image si et seulement si \(\widehat q_k\) est injectif. Pour une variable aléatoire \(Z\) à valeurs dans \(H_k\) et de distribution \(\mu_k\), on a
\[
 \mathsf H(Z)=\mathsf H(q_k(Z))+\mathsf H(Z\mid q_k(Z)).
\]
Si \(\widehat q_k\) est injectif sur le support de \(Z\), alors
\[
 \mathsf H(Z\mid q_k(Z),M_k(Z))=0.
\]
Ici, \(\mathsf H(Z)=-\sum_z\Pr(Z=z)\log\Pr(Z=z)\), avec
\(0\log0=0\), est l'entropie finie en nats ; l'entropie conditionnelle est
la moyenne de l'entropie des distributions conditionnelles.
\end{proposition}

\begin{proof}
Un inverse récupère un unique antécédent, ce qui exige l’injectivité ;
réciproquement, celle-ci permet d’attribuer à chaque paire exprimée son unique
antécédent. Puisque \(q_k(Z)\) est une fonction déterministe de \(Z\),
regrouper la somme qui définit \(\mathsf H(Z)\) par fibres de \(q_k\) donne la
règle de la chaîne. Si la paire \((q_k(Z),M_k(Z))\) distingue le support,
chaque distribution conditionnelle correspondante est concentrée sur un seul
état et son entropie est nulle.

\end{proof}

Ne pas exprimer une coordonnée et effacer le registre qui permet de la récupérer sont
des opérations différentes. Ce résultat n’identifie pas l’entropie à la cardinalité
du continu et n’introduit pas d’interprétation thermique : son domaine, son lecteur et sa
distribution sont ceux qui viennent d’être fixés.


\subsubsection{Équivariance de la mesure sous les symétries de l'arbre}
\label{act21:sec:ch-naturalidad-medida}

\begin{proposition}[Transport conjoint de l'arbre et de sa mesure]
\label{act21:prop:ch-naturalidad-medida}
Soit \(g=(g_k)_{k\geq0}\) une collection de bijections \(g_k:H_k\to H_k\)
qui commute avec les troncatures, fixe la racine et transporte bijectivement
les émissions admissibles, en conservant leurs multiplicités. Alors

\[
 d(g_k\zeta)=d(\zeta),\qquad
 (g_k)_*\mu_k=\mu_k,\qquad (g_\infty)_*\mu=\mu.
\]
Plus généralement, si \(\mu_\zeta\) est la loi des continuations d'une
histoire \(\zeta\in H_k\), on a
\[
 (g_\infty)_*\mu_\zeta=\mu_{g_k\zeta}.
\]
\end{proposition}

\begin{proof}
La bijection des émissions conserve leur nombre avec multiplicités. La
distribution uniforme des germes est également invariante. Chaque trajectoire
et son image reçoivent donc le même produit de poids locaux
\(1/d(\zeta)\). On obtient l'égalité des mesures à chaque niveau et dans
chaque cylindre. La cohérence avec les troncatures définit \(g_\infty\),
et l'unicité de la mesure déterminée par les cylindres prouve l'égalité
à la limite. Le même argument, en commençant en \(\zeta\), prouve
l'affirmation sur les continuations.
\end{proof}

On obtient ainsi la naturalité de l'arbre et de la mesure : un changement cohérent
de représentation transporte les poids et n'oblige pas à les choisir à nouveau.
La conservation des émissions et des multiplicités est une hypothèse essentielle de
la proposition, et non une donnée ajoutée après la construction.
